% Generated draft: requires mathematical and source review.
\documentclass[11pt]{article}
\usepackage[margin=25mm]{geometry}
\usepackage{fontspec}
\setmainfont{Latin Modern Roman}
\usepackage{amsmath,amssymb,mathtools,unicode-math}
\setmathfont{Latin Modern Math}
\usepackage{enumitem,needspace,xurl,hyperref}
\setlength{\parindent}{0pt}
\setlength{\parskip}{5pt}
\setlength{\emergencystretch}{3em}
\newcommand{\N}{\mathbb{N}}
\newcommand{\Z}{\mathbb{Z}}
\newcommand{\Q}{\mathbb{Q}}
\newcommand{\R}{\mathbb{R}}
\newcommand{\C}{\mathbb{C}}
\newcommand{\F}{\mathbb{F}}
\newcommand{\A}{\mathbb{A}}
\newcommand{\PP}{\mathbb{P}}
\newcommand{\E}{\mathbb{E}}
\newcommand{\eps}{\varepsilon}
\newcommand{\dd}{\,\mathrm{d}}
\begin{document}
\section*{UnsolvedMath ID: 30005650}
\textbf{Source problem code:} catalogue entry
\subsection{Definitions and mathematical statement}
For a field \(F\), use the conventional Kronecker dimension
\[
\operatorname{Kdim}(F)=
\begin{cases}
\operatorname{trdeg}_{\Q}(F)+1,&\operatorname{char}(F)=0,\\
\operatorname{trdeg}_{\F_p}(F),&\operatorname{char}(F)=p>0.
\end{cases}
\]
No finite-generation hypothesis is imposed. The supplied excerpt does not explicitly define this dimension convention.

Interpret the source's \emph{inseparable closure} \(F^i\) as its perfect closure: in a fixed algebraic closure \(\overline F\), set
\[
F^i=\begin{cases}
F,&\operatorname{char}(F)=0,\\
\displaystyle\bigcup_{r\geq 0}\{x\in\overline F:x^{p^r}\in F\},&\operatorname{char}(F)=p>0.
\end{cases}
\]
Its abstract field-isomorphism type does not depend on the chosen algebraic closure. This interpretation agrees with the supplied literature assessment's terminology, but the excerpt does not separately define \(F^i\).

Regard the multiplicative abelian group \(F^\times\) as a \(\Z\)-module. Its Milnor \(K\)-theory ring is the graded ring
\[
K_*^M(F)=T_{\Z}(F^\times)\big/\left\langle a\otimes(1-a):a\in F\setminus\{0,1\}\right\rangle,
\]
where the generators \(F^\times\) have degree one and \(K_0^M(F)=\Z\). Write
\[
A(F)=K_*^M(F)\otimes_{\Z}\Q.
\]
An isomorphism of these graded \(\Q\)-algebras is understood to preserve the grading and the unit.

\textbf{Conjecture (reconstruction within the stated class).} For every pair of fields \(F,G\) with \(\operatorname{Kdim}(F)\geq 2\) and \(\operatorname{Kdim}(G)\geq 2\),
\[
A(F)\cong A(G)\quad\text{as graded \(\Q\)-algebras}
\quad\Longrightarrow\quad
F^i\cong G^i\quad\text{as fields}.
\]
This formalizes the source's assertion that the isomorphism type of \(F^i\) is determined by \(A(F)\), interpreting comparison fields as belonging to the same stated class. The source does not explicitly specify whether comparison with fields outside this class is also intended. No compatible lifting of a particular algebra isomorphism, uniqueness, or functorial reconstruction is asserted here.

The accompanying source discussion states that \(F\to F^i\) induces an isomorphism on rationalized Milnor \(K\)-theory. It also discusses indeterminacy from nonzero rational scalings in degree one when \(F^\times\) is divisible; these remarks explain the perfect-closure target and do not add a divisibility hypothesis to the conjecture. The supplied dated assessment reports a reconstruction result under transcendence-degree-at-least-five hypotheses, whose complete additional hypotheses are not supplied, and labels the broader conjecture partially solved. These are supplied contextual claims, not independently verified results in this draft.
\subsection{Short English statement}
The conjecture asks whether rationalized Milnor \(K\)-theory determines the perfect closure of every field of Kronecker dimension at least two, up to abstract field isomorphism. The draft uses conventional meanings of Kronecker dimension and inseparable closure. The supplied partial-solution assessment is dated curation and requires independent review.
\subsection{Sources}
\begin{itemize}
\item Supplied local UnsolvedMath registry record, permanent ID \(30005650\), title \emph{Field Reconstruction from Rationalized Milnor K-Theory}; exact source problem code \texttt{catalogue entry}. The record gives proposed year 2023 and no named proposer.
\item Supplied report citation: \emph{MFO--RIMS Tandem Workshop: Arithmetic Homotopy and Galois Theory (2024)}, Oberwolfach Reports, DOI \texttt{10.4171/owr/2023/42}, pp.\ 2377--2488. \url{https://doi.org/10.4171/owr/2023/42}. The supplied excerpt labels the target Conjecture 1. The citation's 2024 workshop label and its 2023 DOI component are retained as supplied; their relationship has not been checked.
\item Supplied literature source: Topaz, \emph{algebraic dependence and Milnor K-theory}, \url{https://arxiv.org/abs/2211.14665}. The supplied assessment attributes perfect-closure reconstruction under transcendence-degree-at-least-five hypotheses to Topaz, without providing a complete theorem statement.
\item Supplied literature source: OWR report, \url{https://doi.org/10.4171/owr/2023/42}. The exact supplied literature-check date is \texttt{2026-08-22}. The assessment and the label \texttt{partially\_solved} are supplied dated curation, not independently verified here; the record itself asks for independent verification of time-sensitive status and cited claims.
\item The supplied original passage also mentions Bogomolov--Tschinkel [BS08] and Cadoret--Pirutka [CP21], including [CP21, Theorem 4], in background on function-field reconstruction and algebraic dependence. Complete bibliographic entries are absent. The passage describes an initial algebraically closed base-field setting, potentially with additional assumptions, and a subsequent perfect base-field setting; these contextual settings are not hypotheses imposed on Conjecture 1.
\item The registry labels the statement \texttt{corrected\_verified} and reports a check against the official report. This is supplied verification metadata, not a source visit or verification performed for this draft.
\end{itemize}
\end{document}
